% Part: sets-functions-relations
% Chapter: relations
% Section: operations
% Hindi translation (hi-Deva-IN) of Open Logic commit 9620cc73f9c8e0ad003c514a5d3748f29611c4c0.

\documentclass[../../../include/open-logic-section]{subfiles}

\begin{document}

\olfileid[hi]{sfr}{rel}{ops}
\olsection{संबंधों पर संक्रियाएँ}

संबंधों को बदलना या संयोजित करना प्रायः उपयोगी होता है।
\olref[sfr][rel][ord]{prop:stricttopartial} में हमने संबंधों के
\emph{सम्मिलन} पर विचार किया था, जो युग्मों के समुच्चयों के रूप में माने गए
दो संबंधों का सम्मिलन ही है। इसी प्रकार,
\olref[sfr][rel][ord]{prop:partialtostrict} में हमने संबंधों के सापेक्ष अंतर
पर विचार किया था। संबंधों पर की जा सकने वाली कुछ अन्य संक्रियाएँ ये हैं।

\begin{defn}\ollabel{relationoperations} 
मान लीजिए कि $R$, $S$ संबंध हैं और $A$ कोई समुच्चय है।

संबंध $R$ का \emph{प्रतिलोम} $R^{-1} = \Setabs{\tuple{y, x}}{\tuple{x,
    y} \in R}$ है।

संबंधों $R$ और $S$ का \emph{सापेक्ष गुणनफल} $(R \mid S) =
\{\tuple{x, z} : \exists y(Rxy \land Syz)\}$ है।

$R$ का समुच्चय $A$ पर \emph{प्रतिबंधन} $\funrestrictionto{R}{A}= R
\cap A^2$ है।

$R$ का समुच्चय $A$ पर \emph{अनुप्रयोग} $\funimage{R}{A} = \{y :
(\exists x \in A)Rxy\}$ है।
\end{defn}

\begin{ex}
मान लीजिए कि $S \subseteq \Int^2$,~$\Int$ पर उत्तरवर्ती संबंध है, अर्थात
$S = \Setabs{\tuple{x, y} \in \Int^2}{x + 1 = y}$, ताकि $Sxy$ तभी और केवल तभी जब $x + 1 = y$।

$S^{-1}$, $\Int$ पर पूर्ववर्ती संबंध है, अर्थात
$\Setabs{\tuple{x,y}\in\Int^2}{x -1 =y}$।

$S\mid S$ है
$ \Setabs{\tuple{x,y}\in\Int^2}{x + 2 =y}$

$\funrestrictionto{S}{\Nat}$,~$\Nat$ पर उत्तरवर्ती संबंध है।

$\funimage{S}{\{1,2,3\}}$, $\{2, 3, 4\}$ है।
\end{ex}

\begin{defn}[संक्रामक संवरक]मान लीजिए कि $R \subseteq A^2$ एक द्विआधारी संबंध है।
	
संबंध~$R$ का \emph{संक्रामक संवरक} $R^+ = \bigcup_{0 < n \in
\Nat} R^n$ है, जहाँ हम पुनरावर्ती रूप से $R^1 = R$ और $R^{n+1} = R^n
\mid R$ परिभाषित करते हैं।

संबंध $R$ का \emph{स्वतुल्य संक्रामक संवरक} $R^* = R^+ \cup
\Id{A}$ है।
\end{defn}

\begin{ex}
उत्तरवर्ती संबंध $S \subseteq \Int^2$ लीजिए। $S^2xy$ तभी और केवल तभी जब $x + 2 =
y$, $S^3xy$ तभी और केवल तभी जब $x + 3 = y$, इत्यादि। अतः $S^+xy$ तभी और केवल तभी जब
समीकरण $x + n = y$ किसी
$n \geq 1$ के लिए सत्य हो। दूसरे शब्दों में, $S^+xy$ तभी और केवल तभी जब $x < y$, और $S^*xy$ तभी और केवल तभी जब $x \le
y$।
\end{ex}

\begin{prob}
सिद्ध कीजिए कि $R$ का संक्रामक संवरक वास्तव में संक्रामक है।
\end{prob}

\end{document}
